\setlength{\absparsep}{18pt}
\begin{resumo}
\SingleSpacing
{\setlength{\parindent}{1.27cm}\par\noindent\hspace*{\parindent}Este trabalho apresenta o desenvolvimento de uma plataforma experimental de veículo aéreo não tripulado para investigar a integração entre controle de voo e planejamento de missões em ambiente controlado. A arquitetura combina um controlador de voo STM32F411, responsável pela aquisição de sensores, estimação de estados, controle LQG e funções de segurança, com um Raspberry Pi Zero W, responsável pelo planejamento de rotas e pela coordenação sequencial dos comandos de missão. A metodologia compreende modelagem, simulação no MATLAB/Simulink, desenvolvimento embarcado e ensaios de bancada e de voo. Os ensaios dos sensores forneceram dados para a caracterização do ruído de medição, e os testes do conjunto propulsor permitiram caracterizar o empuxo. No ensaio de voo real, observou-se estabilização de rolamento e arfagem; em voo livre, a deriva horizontal e de guinada durante a subida impediu que a referência de altitude fosse alcançada. A comunicação DroneLink foi avaliada pela UART entre o Raspberry Pi e o STM32 reais, com o veículo simulado no microcontrolador e os motores parados: durante os 60 segundos de observação do enlace, não foram registradas perdas nem erros de checksum, e a mediana da latência de ida e volta foi de 10,2 ms. No planejamento de rotas, foram avaliadas 280 instâncias com 3 a 30 pontos. O algoritmo genético obteve o ótimo global em 150 das 200 instâncias em que a solução ótima era conhecida e produziu rotas idênticas no computador pessoal e no Raspberry Pi nas três instâncias comparadas. Os resultados sustentam a viabilidade dos módulos desenvolvidos, enquanto a deriva observada em voo evidencia a necessidade de aperfeiçoar o posicionamento horizontal antes da validação física de missões completas.\par}

\noindent\textbf{Palavras-chave}: veículo aéreo não tripulado. quadricóptero. regulador linear quadrático gaussiano. sistemas embarcados. otimização de rotas.
\end{resumo}

\begin{resumo}[Abstract]
\begin{otherlanguage*}{english}
\SingleSpacing
{\setlength{\parindent}{1.27cm}\par\noindent\hspace*{\parindent}This work presents the development of an experimental unmanned aerial vehicle platform to investigate the integration of flight control and mission planning in a controlled environment. The architecture combines an STM32F411 flight controller, responsible for sensor acquisition, state estimation, LQG control and safety functions, with a Raspberry Pi Zero W, responsible for route planning and the sequential coordination of mission commands. The methodology comprises modelling, MATLAB/Simulink simulation, embedded development, and bench and flight experiments. Sensor experiments provided data for measurement noise characterization, while propulsion tests characterized thrust. The physical flight experiment showed roll and pitch stabilization; in free flight, horizontal and yaw drift during the climb prevented the altitude reference from being reached. DroneLink communication was evaluated over UART between the physical Raspberry Pi and STM32, with the vehicle simulated on the microcontroller and the motors stopped: no packet losses or checksum errors were recorded during the 60-second link observation period, and median round-trip latency was 10.2 ms. Route planning was evaluated on 280 instances with 3 to 30 points. The genetic algorithm found the global optimum in 150 of the 200 instances with a known optimal solution and produced identical routes on the personal computer and Raspberry Pi for the three compared instances. The results support the feasibility of the developed modules, while the drift observed in flight indicates that horizontal positioning requires further improvement before complete missions can be validated in physical flight.\par}

\noindent\textbf{Keywords}: unmanned aerial vehicle. quadrotor. linear quadratic Gaussian regulator. embedded systems. route optimization.
\end{otherlanguage*}
\end{resumo}
